%!TEX root = maindocument.tex

\usepackage[utf8]{inputenc}
\usepackage[ngerman,english]{babel} 
\usepackage[
	backend=bibtex,
	autocite=superscript,
	style=chem-angew,
	articletitle=false,
	biblabel=brackets,
	chaptertitle=false,
	pageranges=false,
	doi=false
	]{biblatex}


\addbibresource{./Literature/Poster.bib}


\usepackage{enumitem}
\usepackage{amsmath,amssymb,nicefrac,braket}
\usepackage{marvosym}
\usepackage{booktabs}
\usepackage{mathtools}
\usepackage{chemmacros}
\usepackage{calc}
\usepackage{scalerel}
\usepackage{lipsum}
\usepackage{lmodern}

\renewcommand{\familydefault}{\sfdefault}


\usepackage{subcaption}
% fix error in subcaption leading to spurious whitespaces
\makeatletter
\renewcommand*\subcaption@label{%
	\caption@withoptargs\subcaption@@label}
\makeatother

\usepackage{xifthen}\usepackage{adjustbox}



\DeclareSIUnit[number-unit-product = {\,}]\Eh{E_\text{h}}

\usetikzlibrary{shapes,arrows,positioning,backgrounds}
\AtBeginEnvironment{tikzfigure}{\captionsetup{type=figure,font=footnotesize}}
\pgfdeclarelayer{background,backgroundlayer}
\pgfsetlayers{background,backgroundlayer,main} 


 \definecolor{RWTHblau100}{RGB}{000,084,159}
 \definecolor{RWTHblau075}{RGB}{064,127,183}
 \definecolor{RWTHblau050}{RGB}{142,186,229}
 \definecolor{RWTHblau025}{RGB}{199,221,242}
 \definecolor{RWTHblau010}{RGB}{232,241,250}
 
 \definecolor{RWTHschwarz100}{RGB}{000,000,000}
 \definecolor{RWTHschwarz075}{RGB}{100,101,103}
 \definecolor{RWTHschwarz050}{RGB}{156,158,159}
 \definecolor{RWTHschwarz025}{RGB}{207,209,210}
 \definecolor{RWTHschwarz010}{RGB}{236,237,237}
 
 \definecolor{RWTHmagenta100}{RGB}{227,000,102}
 \definecolor{RWTHmagenta075}{RGB}{233,096,136}
 \definecolor{RWTHmagenta050}{RGB}{241,158,177}
 \definecolor{RWTHmagenta025}{RGB}{249,210,218}
 \definecolor{RWTHmagenta010}{RGB}{253,250,240}
 
 \definecolor{RWTHgelb100}{RGB}{255,237,000}
 \definecolor{RWTHgelb075}{RGB}{255,240,085}
 \definecolor{RWTHgelb050}{RGB}{255,245,155}
 \definecolor{RWTHgelb025}{RGB}{255,250,209}
 \definecolor{RWTHgelb010}{RGB}{255,253,238}

\definecolor{RWTHpetrol100}{RGB}{000,097,101}
\definecolor{RWTHpetrol075}{RGB}{045,127,131}
\definecolor{RWTHpetrol050}{RGB}{125,164,167}
\definecolor{RWTHpetrol025}{RGB}{191,208,209}
\definecolor{RWTHpetrol010}{RGB}{230,236,236}

\definecolor{RWTHtuerkis100}{RGB}{000,152,161}
\definecolor{RWTHtuerkis075}{RGB}{000,177,183}
\definecolor{RWTHtuerkis050}{RGB}{137,204,207}
\definecolor{RWTHtuerkis025}{RGB}{202,231,231}
\definecolor{RWTHtuerkis010}{RGB}{235,246,246}

\definecolor{RWTHgruen100}{RGB}{087,171,039}
\definecolor{RWTHgruen075}{RGB}{141,192,096}
\definecolor{RWTHgruen050}{RGB}{184,214,152}
\definecolor{RWTHgruen025}{RGB}{221,235,206}
\definecolor{RWTHgruen010}{RGB}{242,247,236}

\definecolor{RWTHmaigruen100}{RGB}{089,205,000}
\definecolor{RWTHmaigruen075}{RGB}{208,217,092}
\definecolor{RWTHmaigruen050}{RGB}{224,230,154}
\definecolor{RWTHmaigruen025}{RGB}{240,243,208}
\definecolor{RWTHmaigruen010}{RGB}{249,250,237}

\definecolor{RWTHorange100}{RGB}{246,168,000}
\definecolor{RWTHorange075}{RGB}{250,190,080}
\definecolor{RWTHorange050}{RGB}{253,212,143}
\definecolor{RWTHorange025}{RGB}{254,234,201}
\definecolor{RWTHorange010}{RGB}{255,247,234}

\definecolor{RWTHrot100}{RGB}{204,007,030}
\definecolor{RWTHrot075}{RGB}{216,092,065}
\definecolor{RWTHrot050}{RGB}{230,150,121}
\definecolor{RWTHrot025}{RGB}{243,205,187}
\definecolor{RWTHrot010}{RGB}{250,235,227}

\definecolor{RWTHbordeaux100}{RGB}{161,016,053}
\definecolor{RWTHbordeaux075}{RGB}{182,082,086}
\definecolor{RWTHbordeaux050}{RGB}{205,139,135}
\definecolor{RWTHbordeaux025}{RGB}{229,197,192}
\definecolor{RWTHbordeaux010}{RGB}{245,232,229}

\definecolor{RWTHviolett100}{RGB}{097,033,088}
\definecolor{RWTHviolett075}{RGB}{131,078,117}
\definecolor{RWTHviolett050}{RGB}{168,133,158}
\definecolor{RWTHviolett025}{RGB}{210,192,205}
\definecolor{RWTHviolett010}{RGB}{237,229,234}

\definecolor{RWTHlila100}{RGB}{122,111,172}
\definecolor{RWTHlila075}{RGB}{155,145,193}
\definecolor{RWTHlila050}{RGB}{188,181,215}
\definecolor{RWTHlila025}{RGB}{222,218,235}
\definecolor{RWTHlila010}{RGB}{242,240,247}

\definecolorstyle{myColorStyle}{
}{
	% Background Colors 
	\colorlet{backgroundcolor}{RWTHschwarz010}%{\white}
	\colorlet{framecolor}{RWTHblau100}
	% Title Colors 
	\colorlet{titlefgcolor}{RWTHschwarz100} 
	\colorlet{titlebgcolor}{white}
	% Block Colors 
	\colorlet{blocktitlebgcolor}{RWTHblau100} 
	\colorlet{blocktitlefgcolor}{white} 
	\colorlet{blockbodybgcolor}{white} 
	\colorlet{blockbodyfgcolor}{RWTHschwarz100}
	% Innerblock Colors 
	\colorlet{innerblocktitlebgcolor}{white} 
	\colorlet{innerblocktitlefgcolor}{RWTHschwarz100} 
	\colorlet{innerblockbodybgcolor}{white} 
	\colorlet{innerblockbodyfgcolor}{RWTHschwarz100}
	% Note colors 
	\colorlet{notefgcolor}{RWTHschwarz100} 
	\colorlet{notebgcolor}{RWTHgelb050} 
	\colorlet{notefrcolor}{RWTHschwarz100}
	%Frame color
	\colorlet{framecolor}{RWTHblau100} 
}

\usecolorstyle{myColorStyle} 


\settitle{\vspace*{1em} \centering %\vbox{
		%\@titlegraphic \\[\TP@titlegraphictotitledistance] \centering
		\color{titlefgcolor} {
			\Huge \sffamily  \bfseries{\@title} \par}
			%\Huge \sffamily \@title \par}
		\vspace*{0.5em}
		{\huge \@author \par} \vspace*{1em} {\LARGE \@institute}
}
%}

\definetitlestyle{titleRWTH}{
	width=1.0\paperwidth, roundedcorners=0, linewidth=0pt, innersep=1cm,
	titletotopverticalspace=0mm, titletoblockverticalspace=20mm,
	titlegraphictotitledistance=10pt
}{
	\draw[
	draw=none,
	line width=0pt,
	%draw, color=RWTHschwarz100, 
	fill=titlebgcolor]%
	(\titleposleft,\titleposbottom) rectangle (\titleposright,\titlepostop); %
}

\usetitlestyle{titleRWTH}


\defineblockstyle{blockRWTH}{
titlewidthscale=1, bodywidthscale=1, titleleft,% titlecenter,
titleoffsetx=0pt, titleoffsety=0pt, bodyoffsetx=0pt, bodyoffsety=2mm,
bodyverticalshift=5mm, roundedcorners=5, linewidth=5pt,
titleinnersep=4mm, bodyinnersep=4mm
}{
\draw[rounded corners=\blockroundedcorners, inner sep=\blockbodyinnersep, line width=\blocklinewidth, color=framecolor, fill=blockbodybgcolor]
(blockbody.south west) rectangle (blockbody.north east); %
\ifBlockHasTitle%
\draw[rounded corners=\blockroundedcorners, inner sep=\blocktitleinnersep, line width=\blocklinewidth, color=framecolor, fill=blocktitlebgcolor]
(blocktitle.south west) rectangle (blocktitle.north east); %
\fi%
}
\useblockstyle{blockRWTH}

\defineinnerblockstyle{innerblockRWTH}{
titlewidthscale=1, bodywidthscale=1, titleleft,% titlecenter,
%titleoffsetx=0pt, titleoffsety=8pt, bodyoffsetx=0pt, bodyoffsety=16pt,
bodyverticalshift=0pt, roundedcorners=0, linewidth=4pt,
titleinnersep=12pt, bodyinnersep=6pt
}{
\begin{scope}[line width=\innerblocklinewidth, rounded
	corners=\innerblockroundedcorners, solid]
	\ifInnerblockHasTitle %
	\draw[color=innerblocktitlebgcolor, fill=innerblocktitlebgcolor]
	(innerblockbody.south west) rectangle (innerblocktitle.north east);
	\draw[color=innerblocktitlebgcolor, fill=innerblockbodybgcolor]
	(innerblockbody.south west) rectangle (innerblockbody.north east);
	\else
	\draw[color=innerblocktitlebgcolor, fill=innerblockbodybgcolor]
	(innerblockbody.south west) rectangle (innerblockbody.north east);
	\fi
\end{scope}
}
\useinnerblockstyle{innerblockRWTH}